% !TeX root = ../main.tex
\section{Amplifier and effect modelling}
\label{sec:modelling}

Modelling is the \emph{forward} problem: given the dry signal, produce the
processed one. It matters here for two reasons: a model can be the output of a
tone estimator (a capture loaded into a player), and a model with few,
interpretable parameters makes the inverse problem easier to search.

\subsection{Black-box models from paired data}

\begin{keyidea}[: capture = learn the device from dry/wet pairs]
Record a few minutes of dry guitar, play it through the real amplifier, record
the output, and train a network to map one to the other with an ESR loss. The
network then imitates that amplifier at that setting. This is what Neural Amp
Modeler, Kemper profiling, TONEX and Neural Capture do.
\end{keyidea}

Feed-forward WaveNets~\cite{damskagg2019deep,damskagg2019realtime} and recurrent
networks~\cite{wright2019realtime} model amplifiers and pedals from a few minutes
of paired audio. Wright et al.\ report test ESR of 0.2--0.3\% and MUSHRA scores
of 90--98 on two tube amplifiers, faster than real time on a CPU
\reported~\cite{wright2020realtime}; a perceptual pre-emphasis improves the
result at no cost~\cite{wright2020perceptual}. Temporal convolutional networks
model compressors~\cite{steinmetz2022efficient}; time-varying feature modulation
extends their memory for fuzz and compression~\cite{comunita2023modelling}. The
largest comparison (40 devices; LSTM, TCN, GCN and S4 models) ranks state-space
models with time-varying modulation best and finds LSTMs unreliable across
devices~\cite{comunita2025differentiable}. Vanhatalo et al.\ review the
field~\cite{vanhatalo2022review}.

Neural Amp Modeler (NAM) has become the open format for captures: an MIT-licensed
trainer, core and plugin~\cite{atkinsonndnam}, a JSON \file{.nam} file with
WaveNet (``A1'', and the slimmable ``A2''~\cite{atkinson2025slimmable,atkinson2026a2})
and LSTM architectures, and optional metadata (gear type, tone type, gain). The
metadata are ready-made weak labels for a capture library. Any capture is one
setting of one device, made with the device present.

\begin{inapp}[: NAM in the app]
The NAM Model pedal runs \file{.nam} files through the pure-Rust \code{nam-rs}
runtime; a standard WaveNet runs at about 11 times real time on one x86 core
and does not resample, so models must match the 48\kHz{} host rate. The tone
estimator does not use NAM captures yet (Section~\ref{sec:design}).
\end{inapp}

\subsection{Knobs and many devices in one model}

A model can take the knobs as inputs. Juvela et al.\ recorded 4.5\,h of an
amplifier with a robot turning its knobs and trained a conditioned LSTM that
listeners could not tell from a circuit simulation~\cite{juvela2023endtoend}.
PANAMA chooses which settings to record by the disagreement of an ensemble, and
with 75 settings matches per-setting NAM captures in a listening
test~\cite{grotschla2025parametric}. Hypernetwork conditioning beats
concatenation for recurrent models~\cite{yeh2024hyper}.

One model can also cover many devices. Chen et al.\ condition a GCN on a tone
embedding from a contrastively trained encoder, so a wet reference clip selects
the tone; on unseen amplifiers the loss is about twice the seen-amplifier loss,
and a case study missed an EQ the model did not
represent~\cite{chen2024zeroshot,chen2024demo}. Open-Amp trains a TCN with a
394-entry device embedding on renders through crowd-sourced captures; a new
analog pedal is added by fitting only its embedding~\cite{wright2025openamp}.
Embedding-only enrolment is a small search space, a candidate target for a
wet-only loss (the paper used paired data).

\subsection{Grey-box and differentiable signal processing}

\term{Grey-box} models keep the structure of the circuit and learn only its
values. The classic block structure is \term{Wiener--Hammerstein}: a linear
filter, a static (memoryless) nonlinearity, another linear
filter~\cite{eichas2017block,eichas2018graybox}, later made
differentiable~\cite{kuznetsov2020differentiable,colonel2022waveshapers}.
Esqueda et al.\ learn the component values of Fender, Marshall and Vox tone
stacks from audio~\cite{esqueda2021differentiable}. Miklánek et al.\ put a
white-box tone stack between a recurrent preamp and power stage; trained on
4--12 minutes covering one to three tone-stack settings, it scored 87--95 MUSHRA
on unseen settings where a black-box network trained on the same data scored
18--47 \reported~\cite{miklanek2023neural}. The DDSP guitar amplifier of Yeh et
al.\ (preamp, tone stack, push--pull power amp, output transformer) needs about
1\,350 operations per sample and beats a recurrent model of the same
cost~\cite{yeh2025ddsp}. Grey-box compressors~\cite{wright2022greybox},
phasers~\cite{carson2023differentiable} and modulation
effects~\cite{carson2026gradient} follow the same idea. Grey-box models are not
yet as accurate as neural networks for distortion and fuzz, though competitive
for compressors~\cite{comunita2025differentiable}.

Libraries make these models reusable: DDSP~\cite{engel2020ddsp}, dasp-pytorch
and DeepAFx-ST~\cite{steinmetz2022style}, GRAFX~\cite{lee2024grafx},
NablAFx~\cite{comunita2025nablafx}, DiffVox~\cite{yu2025diffvox}; Hayes et al.\
review differentiable signal processing~\cite{hayes2024review}. DiffVox fits a
differentiable EQ, compressor, delay and reverb to each of 435 vocal tracks and
studies the distribution of the fitted presets, a recipe that would turn guitar
stems into a prior over settings.

\begin{keyidea}[: why this matters for estimation]
The app's amplifier is itself grey-box-like: gain, a voicing, a three-band tone
stack. That matters less for modelling quality than for estimation: a few
interpretable parameters are what a search can explore and what a player can
adjust afterwards.
\end{keyidea}
